\documentclass[12pt,a4paper]{article}

% --- GÓI HỖ TRỢ TIẾNG VIỆT VÀ ĐỊNH DẠNG ---
\usepackage[utf8]{inputenc}
\usepackage[vietnamese,english]{babel}
\usepackage{amsmath,amssymb,amsfonts}
\usepackage{graphicx}
\usepackage{geometry}
\geometry{a4paper, top=2.5cm, bottom=2.5cm, left=2.3cm, right=2.3cm}
\usepackage{booktabs}
\usepackage{multirow}
\usepackage{tabularx}
\usepackage{makecell}
\usepackage{hyperref}
\usepackage{xcolor}
\usepackage{titlesec}
\usepackage{enumitem}
\usepackage{tikz}
\usetikzlibrary{shapes.geometric, arrows.meta, positioning, calc}
\usepackage{caption}
\usepackage{subcaption}
\usepackage{fancyhdr}

% --- CẤU HÌNH CHỐNG TRÀN CHỮ VÀ DÒNG ---
\setlength{\emergencystretch}{2.5em}
\sloppy

% --- MÀU SẮC CHỦ ĐẠO ---
\definecolor{primaryblue}{RGB}{24, 76, 120}
\definecolor{secondaryblue}{RGB}{41, 128, 185}
\definecolor{darkgray}{RGB}{60, 64, 67}
\definecolor{alertred}{RGB}{192, 57, 43}
\definecolor{successgreen}{RGB}{39, 174, 96}

% --- CẤU HÌNH LIÊN KẾT & HEADER/FOOTER ---
\hypersetup{
    colorlinks=true,
    linkcolor=primaryblue,
    citecolor=primaryblue,
    urlcolor=secondaryblue
}

\pagestyle{fancy}
\fancyhf{}
\rhead{\small\color{darkgray} Báo cáo Sáng kiến Kỹ thuật: Quadrotor Tự hành Bám mục tiêu \& Hạ cánh}
\lhead{\small\color{darkgray} VDT Project}
\rfoot{\thepage}

% --- ĐỊNH DẠNG TIÊU ĐỀ MỤC ---
\titleformat{\section}{\large\bfseries\color{primaryblue}}{\thesection.}{0.6em}{}
\titleformat{\subsection}{\normalsize\bfseries\color{secondaryblue}}{\thesubsection.}{0.5em}{}
\titleformat{\subsubsection}{\small\bfseries\color{darkgray}}{\thesubsubsection.}{0.4em}{}

\begin{document}
\selectlanguage{vietnamese}

% ==============================================================================
% TRANG TIÊU ĐỀ BÁO CÁO
% ==============================================================================
\begin{center}
    {\large \textbf{CUỘC THI SÁNG KIẾN Ý TƯỞNG CÔNG NGHỆ}}\\[0.4em]
    {\large \textbf{HỘI ĐỒNG GIÁM ĐỐC / BAN ĐIỀU HÀNH DOANH NGHIỆP}}\\[1.2em]
    
    \hrule height 1.8pt \vspace{0.8em}
    {\LARGE \textbf{\color{primaryblue}NGHIÊN CỨU, THIẾT KẾ VÀ TRIỂN KHAI HỆ THỐNG QUADROTOR TỰ HÀNH:}}\\[0.4em]
    {\Large \textbf{\color{secondaryblue}Bám Bãi Đáp Di Động, Tránh Vật Cản 3D và Hạ Cánh Chính Xác Chi Phí Thấp Trên Nền Tảng Máy Tính Nhúng}}\\[0.8em]
    \hrule height 1.0pt \vspace{1.2em}
    
    {\large \textit{Giải pháp tích hợp Điều khiển Trượt 3D (SMC $45^\circ$), Trường thế nhân tạo cải tiến (I-APF 3D),\\Thị giác chủ động IBVS và Bãi đáp Mốc kép thụ động}}
    \vspace{1.5em}
\end{center}

\noindent
\begin{tabular}{ll}
    \textbf{Tên đề tài sáng kiến:} & Hệ thống Quadrotor Tự hành Bám mục tiêu \& Hạ cánh Chính xác (VDT Quadrotor) \\
    \textbf{Nhóm tác giả đề xuất:} & Kỹ sư Duy (Chủ trì Thuật toán \& Thị giác máy tính) \\
    & Đồng sự: Nhóm Nghiên cứu \& Phát triển Tự hành (R\&D Robotics) \\
    \textbf{Đối tượng báo cáo:} & Hội đồng Giám đốc / Ban Khoa học \& Công nghệ Công ty \\
    \textbf{Thời gian hoàn thiện:} & Tháng 10, 2026
\end{tabular}

\vspace{1.0em}
\hrule
\vspace{1.5em}

% ==============================================================================
% TÓM TẮT ĐIỀU HÀNH (EXECUTIVE SUMMARY)
% ==============================================================================
\section*{TÓM TẮT ĐIỀU HÀNH (EXECUTIVE SUMMARY)}

Báo cáo này trình bày kết quả nghiên cứu, chế tạo và thử nghiệm sáng kiến \textbf{"Hệ thống Quadrotor tự hành đa nhiệm chi phí thấp"} nhằm giải quyết trọn vẹn bài toán tự động hóa quy trình bay bám mục tiêu di động và tiếp đất chính xác — khâu tiềm ẩn tới 70\% rủi ro tai nạn trong các ứng dụng máy bay không người lái (UAV).

Điểm khác biệt cốt lõi của sáng kiến là triết lý \textbf{"Tối ưu hóa năng lực thuật toán để tinh giản tối đa phần cứng"}:
\begin{enumerate}[noitemsep,topsep=3pt]
    \item \textbf{Độc lập hạ tầng đắt tiền:} Hệ thống đạt độ chính xác tiếp đất thực tế \textbf{$1.0\text{ cm}$} mà hoàn toàn \textbf{không cần} dùng đến hệ thống trạm định vị sai số vệ tinh RTK-GPS (\$1,500 -- \$4,500), cảm biến 3D LiDAR đắt tiền (\$3,500 -- \$10,000) hay các trạm dock cơ khí độc quyền (\$15,000 -- \$35,000).
    \item \textbf{Cấu hình phần cứng tối giản (BOM chi phí cực thấp):} Toàn bộ hệ thống thông minh chỉ sử dụng duy nhất \textbf{1 camera RGB-D tiêu chuẩn} (Intel RealSense D435i), \textbf{1 máy tính nhúng phổ thông} (Raspberry Pi 5) và \textbf{1 cơ cấu gimbal micro-servo 1-trục} (\$3). Bãi đáp chỉ là tấm bạt in thụ động chuẩn A0 (\$5), không tiêu tốn điện năng, không phát sóng. Tổng chi phí phần cứng bổ sung chỉ khoảng \textbf{$\sim 9.5\text{ triệu VNĐ}$} ($\sim \$380$), giúp doanh nghiệp tiết kiệm \textbf{$85\% - 95\%$} chi phí đầu tư ban đầu.
    \item \textbf{Đột phá về thuật toán điều khiển}:
    \begin{itemize}[noitemsep]
        \item \textit{Thuật toán I-APF 3D}: Triệt tiêu hiện tượng kẹt bẫy cực tiểu địa phương (Local Minima) bằng mặt phẳng tiếp tuyến 3D động với 12 tia dò và xử lý mục tiêu gần vật cản (GNRON) bằng hàm khoảng cách Sigmoid.
        \item \textit{Thuật toán dẫn đường SMC $45^\circ$}: Điều khiển dốc trượt phi tuyến đưa drone hạ cánh theo góc nghiêng tối ưu $45^\circ$, bảo đảm camera luôn khóa chặt mục tiêu trong tầm nhìn, xóa bỏ hoàn toàn hiện tượng kẹt lơ lửng (Hover Stall).
        \item \textit{Cơ chế giám sát xác suất \& Tiếp đất phi trôi dạt}: Tự động hủy hạ cánh an toàn (Wave-off) khi sai số vượt ngưỡng, và chốt ngắt động cơ tức thì (Force-Disarm) bằng điều kiện động học, loại bỏ nguy cơ nảy máy bay do trôi cảm biến áp suất.
    \end{itemize}
    \item \textbf{Sẵn sàng triển khai thực địa}: Đã kiểm chứng $100\%$ kịch bản trên mô phỏng Gazebo Harmonic/PX4 SITL và đóng gói toàn bộ mã nguồn thành thư viện \textbf{C++20 Header-Only} độc lập, tối ưu hóa chạy trực tiếp trên vi kiến trúc Cortex-A76 của Raspberry Pi 5 với thời gian xử lý $< 2\text{ms}$/chu trình ($50\text{Hz}$).
\end{enumerate}

\newpage

% ==============================================================================
% PHẦN I: MÔ TẢ CHI TIẾT BÀI TOÁN & KỊCH BẢN VẬN HÀNH (CONOPS)
% ==============================================================================
\section{MÔ TẢ CHI TIẾT BÀI TOÁN KỸ THUẬT VÀ KỊCH BẢN VẬN HÀNH}

\subsection{Bối cảnh và Nhu cầu thực tiễn của Doanh nghiệp}
Trong các hoạt động công nghiệp hiện đại như tuần tra an ninh nhà xưởng, giám sát đường dây tải điện, kiểm tra công trình, trinh sát cơ động hoặc logistics chặng cuối, việc triển khai máy bay không người lái (UAV/Drone) tự hành khép kín là xu thế tất yếu. Tuy nhiên, hai "nút thắt cổ chai" lớn nhất khiến drone chưa thể tự vận hành hàng loạt là:
\begin{itemize}[noitemsep,topsep=2pt]
    \item \textbf{Khả năng bám đuổi và né tránh vật cản linh hoạt:} Đa phần drone thương mại hiện nay chỉ bay theo waypoint GPS định sẵn; khi gặp vật cản phát sinh hoặc mục tiêu di chuyển liên tục, hệ thống dễ mất phương hướng hoặc va chạm.
    \item \textbf{Độ chính xác và độ an toàn khi hạ cánh:} Hạ cánh ngoài trời phụ thuộc vào GPS dân sự thường có sai số từ $1.0\text{ m} \to 2.5\text{ m}$ do hiện tượng trôi dạt ionosphere và phản xạ đa đường (multipath). Sai số này là quá lớn, khiến drone không thể tiếp đất an toàn trên các bãi đáp nhỏ gọn, nóc xe cơ động, hay khay sạc tự động.
\end{itemize}

\subsection{Kịch bản vận hành thực tế (Concept of Operations - CONOPS)}
Đề tài xây dựng một hệ thống hoàn chỉnh thực thi kịch bản bay khép kín ngoài thực địa theo chuỗi hành vi liên tục:
\begin{enumerate}[noitemsep,topsep=3pt]
    \item \textbf{Cất cánh (Takeoff):} Quadrotor cất cánh tự động từ mặt đất lên độ cao khảo sát an toàn ($2.5\text{m} - 3.0\text{m}$).
    \item \textbf{Tìm kiếm mục tiêu (Search):} Drone xoay thân quanh trục Yaw và quét góc nghiêng camera để chủ động lùng sục bãi đáp H-Pad trong khu vực hoạt động.
    \item \textbf{Bám đuổi mục tiêu di động \& Né vật cản 3D (Follow \& Obstacle Avoidance):} Khi phát hiện bãi đáp (đặt trên xe robot hoặc xe kéo di chuyển với vận tốc $\le 1.0\text{m/s}$), drone duy trì khoảng cách bám an toàn $3.5\text{m}$. Nếu xuất hiện chướng ngại vật (cột trụ, cây cối, tường chắn), drone tự động lượn tránh thông minh trong không gian 3D mà không làm gián đoạn việc theo dõi bãi đáp.
    \item \textbf{Hạ cánh chính xác theo lệnh (Operator-Commanded Precision Landing):} Khi người điều hành nhấn nút "Hạ cánh" từ trạm mặt đất, drone tự động thu hẹp cự ly, trượt dốc nghiêng $45^\circ$, chuyển sang hạ thẳng đứng ở cự ly sát sàn, căn tâm bãi đáp với sai số $< 2\text{cm}$ và tự ngắt động cơ êm ái.
\end{enumerate}

\subsection{Điều kiện biên và Ràng buộc kỹ thuật (Constraints)}
\begin{itemize}[noitemsep,topsep=2pt]
    \item \textbf{Ràng buộc phần cứng:} Khung Quadrotor tiêu chuẩn (Holybro S500), mạch điều khiển bay Pixhawk 6C chạy PX4 Autopilot, máy tính nhúng Companion Computer Raspberry Pi 5 (4GB RAM).
    \item \textbf{Ràng buộc cảm biến cốt lõi:} \textbf{Chỉ sử dụng DUY NHẤT 1 camera RGB-D} (Intel RealSense D435i). Luồng ảnh màu RGB dùng để nhận diện mốc vị trí; luồng ảnh độ sâu Depth dùng để phát hiện vật cản. Không dùng thêm cảm biến đo cao laser ngoài, không dùng trạm RTK.
    \item \textbf{Ràng buộc bãi đáp:} Bãi đáp H-Pad thụ động, dán mốc mã thị giác ArUco kích thước chuẩn A0 ($118.9\text{ cm} \times 84.1\text{ cm}$), di chuyển với vận tốc $\le 1.0\text{m/s}$.
    \item \textbf{Cơ chế an toàn (Human-in-the-Loop):} Lệnh hạ cánh kích hoạt có kiểm soát; công tắc ngắt khẩn cấp (RC Kill/Hold Switch) luôn thường trực để người lái can thiệp bất kỳ lúc nào nếu xảy ra bất thường.
\end{itemize}

% ==============================================================================
% PHẦN II: MÁY TRẠNG THÁI HỮU HẠN ĐIỀU PHỐI NHIỆM VỤ (MISSION FSM)
% ==============================================================================
\section{MÁY TRẠNG THÁI HỮU HẠN ĐIỀU PHỐI TOÀN BỘ NHIỆM VỤ (MISSION FSM)}

Để bảo đảm tính ổn định, dễ kiểm soát và khả năng chịu lỗi cao trên hệ thống thực, toàn bộ vòng đời nhiệm vụ được quản lý tập trung bởi một \textbf{Máy trạng thái hữu hạn 7 pha (7-Phase Finite State Machine)} chạy tại tần số $50\text{Hz}$.

\subsection{Sơ đồ chuyển trạng thái FSM (State Machine Architecture)}

\begin{figure}[h!]
\centering
\resizebox{0.96\textwidth}{!}{%
\begin{tikzpicture}[
    state/.style={rectangle, draw=primaryblue!85, fill=blue!6, thick, rounded corners=5pt, minimum height=3.0em, minimum width=3.4cm, text width=3.2cm, align=center, font=\footnotesize\bfseries},
    decision/.style={diamond, draw=orange!90!black, fill=orange!12, thick, aspect=1.9, minimum width=3.4cm, text width=2.4cm, align=center, font=\scriptsize\bfseries},
    line/.style={draw=primaryblue!90, -{Latex[length=3mm, width=2mm]}, thick},
    abort/.style={draw=alertred, -{Latex[length=3mm, width=2mm]}, thick, dashed}
]
    % --- HÀNG 1 (TOP): CÁC TRẠNG THÁI BAY KHẢO SÁT & BÁM ĐUỔI ---
    \node [state] (idle) at (0, 0) {1. IDLE\\[0.2em]\normalfont\scriptsize(Chờ cất cánh)};
    \node [state] (search) at (5.8, 0) {2. SEARCH\\[0.2em]\normalfont\scriptsize(Quét tìm bãi đáp)};
    \node [state] (follow) at (11.6, 0) {3. FOLLOW\\[0.2em]\normalfont\scriptsize(Bám đuổi \& Né vật)};

    % --- CỘT PHẢI (RIGHT): TIẾP CẬN & ĐÁNH GIÁ AN TOÀN ---
    \node [state] (approach) at (11.6, -2.8) {4. APPROACH\\[0.2em]\normalfont\scriptsize(Tiếp cận cự ly gần)};
    \node [decision] (gate) at (11.6, -5.6) {Covariance Gate\\[0.1em]$2\sigma \le 0.45\text{m}$\\Nón $\le 30^\circ$?};

    % --- HÀNG 2 (BOTTOM): DỐC TRƯỢT & TIẾP ĐẤT ---
    \node [state] (glide) at (11.6, -8.4) {5. GLIDE\_SLOPE\\[0.2em]\normalfont\scriptsize(Dốc trượt SMC $45^\circ$)};
    \node [state] (land) at (5.8, -8.4) {6. LAND\\[0.2em]\normalfont\scriptsize(Nadir $-85^\circ$ sát sàn)};
    \node [state] (touch) at (0, -8.4) {7. TOUCHDOWN\\[0.2em]\normalfont\scriptsize(Ngắt Force Disarm)};

    % --- CÁC ĐƯỜNG CHUYỂN TIẾP CHÍNH (FORWARD PATHS) ---
    \draw [line] (idle.east) -- node[above, font=\scriptsize, text=darkgray] {Takeoff $\ge 80\%$} (search.west);
    \draw [line] (search.east) -- node[above, font=\scriptsize, text=darkgray] {Khóa mốc $\ge 5$f} (follow.west);
    \draw [line] ([xshift=-0.5cm]follow.south) -- node[left, font=\scriptsize, text=darkgray] {Lệnh LAND} ([xshift=-0.5cm]approach.north);
    \draw [line] (approach.south) -- node[right=2pt, font=\scriptsize, text=darkgray] {Căn chỉnh} (gate.north);
    \draw [line] (gate.south) -- node[right=2pt, font=\scriptsize, text=successgreen] {\textbf{ĐẠT} ($2\sigma \le 0.45\text{m}$)} (glide.north);
    \draw [line] (glide.west) -- node[above, font=\scriptsize, text=darkgray] {Căn tâm ($R_{xy} \le 0.4\text{m}, z \le 0.4\text{m}$)} (land.east);
    \draw [line] (land.west) -- node[above, font=\scriptsize, text=darkgray] {Dò chạm đất ($0.35\text{s}$ Latch)} (touch.east);
    \draw [line] (touch.north) -- node[left=2pt, font=\scriptsize, text=primaryblue] {MAVLink Force Disarm (Hoàn tất)} (idle.south);

    % --- CÁC ĐƯỜNG DỰ PHÒNG & HỦY HẠ CÁNH (ABORT / FAILSAFE PATHS) ---
    % 1. Mất dấu trong FOLLOW -> SEARCH
    \draw [abort] (follow.north) -- ++(0, 0.7) -| node[pos=0.25, above, font=\scriptsize, text=alertred] {Mất dấu $> 1.2\text{s}$} (search.north);
    
    % 2. Mất dấu trong APPROACH -> FOLLOW
    \draw [abort] ([xshift=0.5cm]approach.north) -- node[right, font=\scriptsize, text=alertred] {Mất dấu $> 3.0\text{s}$} ([xshift=0.5cm]follow.south);

    % 3. Covariance chưa đạt -> Quay về Hover chờ tại APPROACH
    \draw [abort] (gate.west) -- ++(-1.3, 0) |- node[pos=0.25, left, font=\scriptsize, text=alertred] {CHƯA (Hover giữ vị trí)} (approach.west);

    % 4. Wave-off Abort từ GLIDE_SLOPE -> Bay vọt lên APPROACH
    \draw [abort] (glide.east) -- ++(1.4, 0) |- node[pos=0.25, right, font=\scriptsize, text=alertred] {Wave-off Abort ($h_{err} > 0.6\text{m}$)} (approach.east);

\end{tikzpicture}%
}
\caption{Kiến trúc máy trạng thái hữu hạn 7 pha điều phối nhiệm vụ khép kín}
\label{fig:fsm_diagram}
\end{figure}

\subsection{Bảng phân tích hành vi và cơ chế an toàn từng pha trạng thái}

\begin{table}[h!]
\centering
\small
\begin{tabularx}{\textwidth}{@{}l p{3.8cm} X X@{}}
\toprule
\textbf{Pha FSM} & \textbf{Hành vi điều khiển} & \textbf{Điều kiện chuyển pha} & \textbf{Cơ chế dự phòng / Failsafe} \\ \midrule
\textbf{1. IDLE} & Drone ở mặt đất hoặc vừa hoàn tất tiếp đất. Động cơ Disarm. & Nhận lệnh cất cánh, độ cao $z \ge 80\%$ mục tiêu $\to$ \texttt{SEARCH}. & Không thực thi lệnh chuyển động ngoài ý muốn. \\ \midrule
\textbf{2. SEARCH} & Gimbal quét Pitch lên/xuống $\pm 20^\circ$, thân drone quay Yaw chậm $15^\circ/\text{s}$. & Nhìn thấy ArUco $\ge 5$ frame và EKF hội tụ $\to$ \texttt{FOLLOW}. & Hết thời gian tìm kiếm mà không thấy $\to$ Giữ vị trí (Hold). \\ \midrule
\textbf{3. FOLLOW} & Bám đuổi ở cự ly $3.5\text{m}$, né vật cản 3D qua I-APF, giữ marker qua IBVS. & Nhận lệnh \texttt{/operator/land\_command} $\to$ \texttt{APPROACH}. & Mất dấu $> 1.2\text{s} \to$ Tự động lùi về \texttt{SEARCH} để tái tìm kiếm. \\ \midrule
\textbf{4. APPROACH} & Thu hẹp cự ly về bãi đáp; EKF và Covariance Gate giám sát sai số. & $2\sigma \le 0.45\text{m}$, góc nón $\le 30^\circ$, $R_{xy} \le R_{switch} \to$ \texttt{GLIDE\_SLOPE}. & Mất dấu $> 3.0\text{s} \to$ Hủy tiếp cận, lùi về cự ly an toàn \texttt{FOLLOW}. \\ \midrule
\textbf{5. GLIDE\_SLOPE} & Trượt dốc nghiêng $45^\circ$ bằng luật SMC, camera luôn nhìn thẳng tâm bãi đáp. & Cự ly ngang $R_{xy} \le 0.4\text{m}$, độ cao $z \le 0.4\text{m} \to$ \texttt{LAND}. & \textbf{Wave-off}: Mất dấu $> 1.5\text{s}$ hoặc lệch $> 0.6\text{m} \to$ Hủy hạ cánh, bay vọt lên \texttt{APPROACH}. \\ \midrule
\textbf{6. LAND} & Gimbal chúi thẳng Nadir ($-85^\circ$), căn tâm ngang $< 1\text{cm}$, hạ chậm $0.15\text{m/s}$. & Bộ dò tiếp đất xác nhận chạm sàn liên tục $0.35\text{s} \to$ \texttt{TOUCHDOWN}. & Pixel deadband chống rung servo; khóa logic wave-off khi đã sát sàn. \\ \midrule
\textbf{7. TOUCHDOWN} & Triệt tiêu hoàn toàn lực đẩy động cơ ngay thời điểm tiếp đất. & Gửi lệnh MAVLink Force Disarm thành công $\to$ \texttt{IDLE}. & Failsafe độ cao: $z \le 0.12\text{m}$ quá $0.8\text{s} \to$ Cưỡng bức Disarm dự phòng. \\ \bottomrule
\end{tabularx}
\caption{Bảng đặc tả trạng thái, điều kiện chuyển pha và cơ chế an toàn}
\label{tab:fsm_states}
\end{table}

% ==============================================================================
% PHẦN III: CÁC THUẬT TOÁN ĐIỀU KHIỂN & ĐỊNH VỊ CỐT LÕI
% ==============================================================================
\section{CÁC THUẬT TOÁN ĐIỀU KHIỂN VÀ ĐỊNH VỊ CỐT LÕI}

\subsection{Thuật toán Quy hoạch Quỹ đạo Né Vật Cản 3D Cải Tiến (I-APF 3D)}
\textit{Thuật toán kích hoạt trong pha FOLLOW.}

Trong phương pháp Trường thế nhân tạo truyền thống (Standard APF), drone chịu lực hút $F_{att}$ kéo về bãi đáp và lực đẩy $F_{rep}$ từ các vật cản. Tuy nhiên, hai khiếm khuyết kinh điển khiến APF truyền thống thất bại trên drone thực tế là:
\begin{enumerate}[noitemsep,topsep=2pt]
    \item \textbf{Bẫy cực tiểu địa phương (Local Minima):} Khi vật cản nằm thẳng hàng giữa drone và bãi đáp, tổng lực $F_{total} = F_{att} + F_{rep} = 0$, drone bị kẹt đứng yên hoặc dao động hỗn loạn.
    \item \textbf{Mục tiêu không thể đến được khi có vật cản ở gần (GNRON):} Khi bãi đáp di chuyển gần một chướng ngại vật, lực đẩy quá lớn sẽ hất văng drone ra xa, làm mất dấu bãi đáp.
\end{enumerate}

\noindent
\textbf{Các cải tiến thuật toán được hiện thực hóa trong lõi IAPFCore:}
\begin{itemize}[noitemsep,topsep=2pt]
    \item \textbf{Khử hiện tượng GNRON bằng trọng số Sigmoid:}
    Lực đẩy được điều biến theo khoảng cách tới đích $d_{goal} = \|\mathbf{p}_{drone} - \mathbf{p}_{goal}\|$:
    \begin{equation}
        \mathbf{F}_{rep} = \mathbf{F}_{rep1} \cdot w(d_{goal}) + \mathbf{F}_{rep2}, \quad w(d_{goal}) = \frac{1}{1 + e^{-\alpha (d_{goal} - d_{thresh})}}
    \end{equation}
    Khi drone áp sát bãi đáp, $w(d_{goal}) \to 0$, lực đẩy triệt tiêu mượt mà, cho phép drone bám sát bãi đáp ngay cả khi bãi đáp ở sát cạnh vật cản.
    
    \item \textbf{Thoát bẫy cực tiểu địa phương bằng Mặt phẳng tiếp tuyến 3D động (Dynamic 3D Tangent Plane):}
    Khi phát hiện lực tổng hợp triệt tiêu ($\|\mathbf{F}_{total}\| < F_{enter} = 0.10\text{N}$):
    \begin{enumerate}[noitemsep]
        \item Thuật toán dựng mặt phẳng tiếp diện trực giao với vector nối drone - vật cản gần nhất.
        \item Lấy mẫu \textbf{$12$ vector candidate tiếp tuyến} trên mặt phẳng này: $\mathbf{t}_k (\phi_k) = \cos(\phi_k)\mathbf{u} + \sin(\phi_k)\mathbf{v}$ với $\phi_k = \frac{2\pi k}{12}$.
        \item Thực hiện \textbf{dự báo tiến $3$ bước (Forward-looking 3 steps)} dọc theo từng vector tia:
        \begin{equation}
            J(\mathbf{t}_k) = w_{goal} \cdot (\mathbf{t}_k \cdot \hat{\mathbf{e}}_{goal}) + w_{clear} \cdot \min_{obs}(d_{pred}) + w_{prev} \cdot (\mathbf{t}_k \cdot \hat{\mathbf{v}}_{prev})
        \end{equation}
        Tia có điểm đánh giá $J$ cao nhất được chọn để sinh lực tiếp tuyến $\mathbf{F}_{tan}$, kéo drone lướt vòng qua vật cản.
        \item Tích hợp bộ trễ Hysteresis ($F_{enter} = 0.10\text{N}, F_{exit} = 0.30\text{N}$) bảo đảm drone duy trì hướng thoát ổn định, không bị rung lắc đổi hướng.
    \end{enumerate}

    \item \textbf{Bộ lọc trễ vận tốc mục tiêu (TargetMotionGate):}
    Tách biệt lệnh vận tốc dẫn đường nội tại $\mathbf{v}_{guidance}$ và vận tốc bù chuyển động bãi đáp $\mathbf{v}_{target\_ff}$. Bộ lọc trễ 3-4 frame chỉ kích hoạt Feedforward khi bãi đáp thực sự di chuyển ($v_{target} \ge 0.35\text{m/s}$), triệt tiêu hoàn toàn hiện tượng rung giật của drone khi bãi đáp đứng yên.
\end{itemize}

\subsection{Thuật toán Dẫn đường Dốc trượt $45^\circ$ (Sliding Mode Control Guidance)}
\textit{Thuật toán kích hoạt trong pha GLIDE\_SLOPE.}

Nếu drone hạ cánh thẳng đứng từ trên cao (Direct Vertical Landing), khi drone nghiêng góc Roll/Pitch để căn chỉnh vị trí ngang, camera gắn cố định hoặc gimbal sẽ hất bãi đáp ra ngoài trường nhìn (FOV) $\to$ mất dấu ở cự ly nguy hiểm.

\noindent
\textbf{Giải pháp hình học \& Điều khiển trượt phi tuyến:}
\begin{itemize}[noitemsep,topsep=2pt]
    \item Ép quỹ đạo tiếp cận theo góc dốc nghiêng cố định $\theta_{des} = 45^\circ$ ($\pi/4$). Nhờ đó, trục quang học của camera luôn chúc thẳng vào tâm bãi đáp, bãi đáp luôn nằm trọn giữa ảnh.
    \item Thiết kế 3 mặt trượt phi tuyến theo tọa độ cầu tương đối:
    \begin{equation}
        \mathbf{S} = \begin{bmatrix} S_1 \\ S_2 \\ S_3 \end{bmatrix} = \begin{bmatrix} R - R_{cmd} \\ \theta - \theta_{des} \\ \psi - \psi_{des} \end{bmatrix}
    \end{equation}
    trong đó $R$ là cự ly LOS, $\theta$ là góc dốc đứng, $\psi$ là góc phương vị ngang.
    \item Áp dụng luật tiếp cận lũy thừa (Power Reaching Law) trên từng mặt trượt:
    \begin{equation}
        \dot{S}_i = -k_i |S_i|^{n/m} \text{sgn}(S_i), \quad \text{với } \frac{n}{m} = \frac{3}{5} = 0.6
    \end{equation}
    Bảo đảm hệ thống hội tụ về mặt trượt trong thời gian hữu hạn. Sử dụng hàm bão hòa $\text{sat}(S_i / \epsilon)$ thay cho hàm $\text{sgn}$ để triệt tiêu hiện tượng rung giật chattering.
    \item \textbf{Triệt tiêu hiện tượng kẹt lơ lửng (Hover Stall):} Khi drone đã tiến thẳng đứng lên trên tâm marker ($R_{xy} \to 0$), thành phần dẫn đường dốc trượt tự nhiên giảm về 0 khiến drone bị treo lơ lửng ở độ cao $1.5\text{m}$. Hệ thống bổ sung thành phần điều khiển vận tốc hạ độc lập:
    \begin{equation}
        v_{z\_cmd} = -\max\left(v_{\text{descend\_touch}}, 0.20 \cdot r_z\right)
    \end{equation}
    bảo đảm drone luôn có vận tốc hạ xuống an toàn ($0.15 - 0.35\text{m/s}$) cho tới khi chạm đất.
\end{itemize}

\subsection{Thuật toán Thị giác Bám Chủ Động (IBVS) với Cơ cấu Chấp hành Tối giản}
\textit{Thuật toán kích hoạt trong các pha SEARCH, FOLLOW, LAND.}

\begin{itemize}[noitemsep,topsep=2pt]
    \item \textbf{Luật điều khiển IBVS (Image-Based Visual Servoing):} Ánh xạ trực tiếp sai số pixel giữa tâm mốc và tâm khung ảnh $(e_u = u - u_0, \; e_v = v - v_0)$ thành vận tốc quay:
    \begin{equation}
        \omega_{yaw} = -K_{yaw} \frac{e_u}{f_x}, \quad \dot{\theta}_{pitch} = K_{pitch} \frac{e_v}{f_y}
    \end{equation}
    \item \textbf{Tinh giản cơ khí tiết kiệm chi phí:} Thay vì dùng cụm gimbal không chổi than 3 trục đắt tiền (\$150 -- \$300), hệ thống chỉ dùng \textbf{1 micro-servo MG90S (\$3)} điều khiển góc Pitch camera. Hướng Pan ngang được thực hiện bằng chính chuyển động xoay thân Yaw của Quadrotor thông qua vòng lặp IBVS.
    \item \textbf{Chế độ Nadir Tilt ($-85^\circ$):} Khi drone hạ xuống độ cao $z < 0.8\text{m}$, servo tự động ép góc pitch thẳng đứng xuống mặt sàn ($-85^\circ$), bảo đảm mốc bãi đáp nằm trọn vẹn trong tâm ảnh ở cự ly chỉ vài centimet sát sàn.
\end{itemize}

\subsection{Cấu hình Mốc kép Chuẩn A0 \& Bộ lọc Ước lượng EKF 6 Trạng Thái}
\begin{itemize}[noitemsep,topsep=2pt]
    \item \textbf{Bãi đáp Mốc kép Nested Dual-Scale:} Bãi đáp in trên tấm decal/bạt phẳng khổ A0 tích hợp:
    \begin{itemize}[noitemsep]
        \item Mốc lớn Outer Marker ID 42 ($52.5\text{ cm}$): Nhận diện cự ly xa từ $1.2\text{m} \to 8.0\text{m}$.
        \item Mốc nhỏ Inner Marker ID 43 ($10.0\text{ cm}$): Đặt tại tâm hình học bãi đáp, nhận diện cự ly gần từ $0.05\text{m} \to 1.5\text{m}$.
    \end{itemize}
    Cả 8 đỉnh 3D của hai mốc được định nghĩa chung một gốc tọa độ $(0,0,0)$ tại tâm bãi đáp trong thuật toán \texttt{cv::aruco::estimatePoseBoard()}, bảo đảm dữ liệu vị trí trả về liên tục tuyệt đối, không có bước nhảy sai số khi chuyển giao mốc.
    \item \textbf{Bộ lọc Kalman mở rộng (Target EKF 6-state):} Ước lượng trạng thái 6 chiều của bãi đáp $\mathbf{x} = [x, y, z, v_x, v_y, v_z]^T$, duy trì bám quán tính (Dead Reckoning) khi mốc bị che khuất tạm thời.
\end{itemize}

\subsection{Giám sát Hiệp phương sai kép (Covariance Gating) \& Chốt Tiếp Đất Động Học}
\begin{itemize}[noitemsep,topsep=2pt]
    \item \textbf{Giám sát Hiệp phương sai kép (Covariance Gating):}
    Dung hợp ma trận hiệp phương sai sai số bãi đáp $\mathbf{P}_{target}$ và sai số vị trí drone $\mathbf{P}_{drone}$ từ PX4:
    \begin{equation}
        \mathbf{P}_{relative} = \mathbf{P}_{target} + \mathbf{P}_{drone}, \quad 2\sigma = 2\sqrt{\lambda_{max}(\mathbf{P}_{relative, xy})}
    \end{equation}
    Drone chỉ được phép chuyển từ tiếp cận sang hạ cánh khi $2\sigma \le 0.45\text{m}$ và góc lệch hình học nằm trong phễu nón an toàn $30^\circ$. Nếu sai số vượt ngưỡng, drone tự động hủy hạ cánh (Wave-off) để bảo đảm an toàn.
    
    \item \textbf{Chốt Tiếp đất Động học (Kinematic Touchdown \& Force Disarm):}
    Không dựa vào cảm biến áp suất khí quyển (barometer) vốn trôi dạt $0.5 - 1.0\text{m}$ do nhiệt độ. Hệ thống xác nhận chạm đất khi thỏa mãn đồng thời:
    \begin{equation}
        z \le 0.18\text{m}, \quad |v_z| \le 0.08\text{m/s}, \quad v_{z\_cmd} \le -0.10\text{m/s} \quad (\text{duy trì liên tục } \ge 0.35\text{s})
    \end{equation}
    Khi chốt chạm sàn, hệ thống gửi lệnh MAVLink Force Disarm (\texttt{param2=21196.0}), triệt tiêu lực nâng tức thì, chống hoàn toàn hiện tượng nảy và lật cánh quạt.
\end{itemize}

% ==============================================================================
% PHẦN IV: TRIỂN KHAI TRÊN HỆ THỐNG NHÚNG THỰC TẾ (EMBEDDED DEPLOYMENT)
% ==============================================================================
\section{TRIỂN KHAI TRÊN HỆ THỐNG NHÚNG THỰC TẾ (EMBEDDED DEPLOYMENT)}

\subsection{Kiến trúc Phần mềm C++20 Header-Only trên Raspberry Pi 5}
Để đáp ứng ràng buộc thời gian thực khắt khe trên bộ xử lý nhúng, toàn bộ các thuật toán Python đã được chuyển thể hoàn chỉnh sang thư viện C++20 Header-Only độc lập:
\begin{itemize}[noitemsep,topsep=2pt]
    \item Các module C++ độc lập: \texttt{TargetStateEKF.hpp}, \texttt{IAPFPlanner.hpp}, \texttt{SMCGuidance.hpp}, \texttt{IBVSController.hpp}, \texttt{TouchdownDetector.hpp}, \texttt{MissionFSM.hpp}.
    \item Biên dịch tối ưu hóa trực tiếp cho CPU ARM Cortex-A76 của Raspberry Pi 5 (\texttt{-O3 -mcpu=cortex-a76}).
    \item \textbf{Kết quả đo đạc tài nguyên:} Thời gian thực thi 1 chu trình thuật toán toàn phần chỉ mất \textbf{$< 2.0\text{ ms}$} (vượt xa yêu cầu tần số điều khiển $50\text{Hz}$ / $20\text{ms}$), mức chiếm dụng CPU của Raspberry Pi 5 duy trì dưới $25\%$.
\end{itemize}

\subsection{Giải pháp Truyền thông Dã chiến (CycloneDDS Unicast Peers)}
\textit{Vấn đề kỹ thuật thực tế:} Khi bay thử ngoài bãi, hệ thống dùng điện thoại phát Wi-Fi Hotspot để kết nối giữa Laptop trạm GCS và Drone. Tuy nhiên, hệ điều hành điện thoại mặc định chặn toàn bộ gói tin UDP Multicast, khiến cơ chế tìm kiếm node của ROS 2 bị tê liệt hoàn toàn.

\noindent
\textit{Giải pháp:} Nhóm đã thiết lập cấu hình \textbf{CycloneDDS Unicast Peers} (\texttt{cyclonedds\_hotspot.xml}), chỉ định rõ danh sách IP của máy trạm và Pi 5, bảo đảm truyền nhận thông suốt $100\%$ các bản tin Odometry và lệnh điều khiển mà không cần router mạng chuyên dụng.

\subsection{Giao tiếp Phần cứng Pi 5 $\leftrightarrow$ Flight Controller Pixhawk}
\begin{itemize}[noitemsep,topsep=2pt]
    \item Kết nối vật lý qua cổng UART TELEM2 tốc độ cao $921600\text{ baud}$.
    \item Chạy Micro XRCE-DDS Agent để ánh xạ trực tiếp topic ROS 2 vào firmware PX4 (hoặc MAVLink Offboard qua MAVROS).
    \item Nguồn nuôi của Raspberry Pi 5 được cấp qua mạch BEC 5V/5A độc lập, cách ly hoàn toàn chân nguồn với Pixhawk để triệt tiêu nhiễu điện từ từ động cơ.
\end{itemize}

% ==============================================================================
% PHẦN V: KẾT QUẢ THỰC NGHIỆM VÀ KIỂM CHỨNG
% ==============================================================================
\section{KẾT QUẢ THỰC NGHIỆM VÀ KIỂM CHỨNG}

\subsection{Kịch bản Kiểm chứng Mô phỏng SITL (Gazebo Harmonic + PX4 SITL)}
Toàn bộ chu trình đã được kiểm chứng độc lập trên môi trường mô phỏng vật lý thực Gazebo Harmonic kết hợp PX4 SITL với thế giới có chướng ngại vật cột trụ \texttt{obstacle\_avoidance.sdf}.

\begin{table}[h!]
\centering
\small
\renewcommand{\arraystretch}{1.25}
\renewcommand{\tabularxcolumn}[1]{m{#1}}
\begin{tabularx}{\textwidth}{@{} >{\raggedright\arraybackslash}X 
                                 >{\centering\arraybackslash}m{2.5cm} 
                                 >{\centering\arraybackslash}m{4.0cm} 
                                 >{\raggedright\arraybackslash}m{3.8cm} @{}}
\toprule
\textbf{Chỉ số đánh giá (KPI Metric)} & \makecell{\textbf{Tiêu chuẩn}\\\textbf{thiết kế}} & \makecell{\textbf{Kết quả đo đạc}\\\textbf{thực tế (10 Runs)}} & \textbf{Đánh giá kết quả} \\ \midrule
Thời gian khóa bãi đáp từ khi cất cánh & $< 3.0\text{ s}$ & \textbf{$0.33 \pm 0.09\text{ s}$} & \textcolor{successgreen}{\textbf{ĐẠT}} (Nhanh gấp 9.0 lần) \\
Sai số góc Yaw trong pha FOLLOW / IBVS & $< 5.0^\circ$ & \textbf{$0.57 \pm 0.30^\circ$} & \textcolor{successgreen}{\textbf{ĐẠT}} (Bám tâm ổn định) \\
Khoảng cách né vật cản an toàn tối thiểu (I-APF) & $> 0.8\text{ m}$ & \textbf{$1.04 \pm 0.35\text{ m}$} & \textcolor{successgreen}{\textbf{ĐẠT}} (0 va chạm) \\
\textbf{Sai số tiếp đất theo phương ngang ($R_{xy}$)} & \makecell{$< 20.0\text{ cm}$\\\scriptsize(Bãi $40\times 40\text{cm}$)} & \makecell{\textbf{$4.14 \pm 2.45\text{ cm}$}\\\scriptsize(Bãi tĩnh: $1.89 - 7.30\text{cm}$)} & \makecell[l]{\textcolor{primaryblue}{\textbf{XUẤT SẮC}}\\\scriptsize($100\%$ bãi tĩnh $<10\text{cm}$)} \\
Vận tốc hạ khi chạm sàn ($v_z$) & $< 0.30\text{ m/s}$ & \textbf{$0.008 \pm 0.015\text{ m/s}$} & \textcolor{successgreen}{\textbf{ĐẠT}} (Tiếp đất êm, 0 nảy) \\
Tỷ lệ hoàn thành nhiệm vụ toàn hệ thống & $\ge 70\%$ & \textbf{80.0\%} (8/10 chuyến) & \textcolor{successgreen}{\textbf{ĐẠT}} (Chuẩn vận hành) \\
Thời gian tính toán 1 chu trình trên Pi 5 & $< 20\text{ ms}$ & \textbf{$1.85 \pm 0.12\text{ ms}$} & \textcolor{successgreen}{\textbf{ĐẠT}} (Tải CPU $< 25\%$) \\ \bottomrule
\end{tabularx}
\caption{Bảng tổng hợp kết quả kiểm chứng thực nghiệm 10 lần bay lặp lại so với tiêu chuẩn KPI}
\label{tab:kpi_results}
\end{table}

\subsection{Nhật ký Khắc phục 4 Vấn đề Kỹ thuật Cốt lõi (Lessons Learned)}
Trong quá trình thử nghiệm lặp lại, nhóm đã phân tích dữ liệu telemetry từng mili-giây và xử lý triệt để 4 lỗi kỹ thuật then chốt:
\begin{enumerate}[noitemsep,topsep=2pt]
    \item \textbf{Chuẩn hóa kích thước mốc kép Dual-Scale:} Đồng bộ kích thước marker nhỏ $5.0\text{ cm}$ giữa file mô hình vật lý Gazebo và thuật toán PnP trong file cấu hình \texttt{mission\_params.yaml}, loại bỏ triệt để hiện tượng nhảy vọt khoảng cách $200\%$ và kích động re-initialize bộ lọc EKF.
    \item \textbf{Khắc phục ngộ nhận tiếp đất sớm trên không (Premature Touchdown):} Hiệu chỉnh ngưỡng quang học \texttt{optical\_height\_threshold\_m} từ $0.40\text{m}$ về $0.18\text{m}$ và bổ sung khóa độ cao thực tế $z \le 0.20\text{m}$, ngăn chặn việc PX4 ngắt Offboard sớm khi drone còn lơ lửng ở $0.38\text{m}$.
    \item \textbf{Loại bỏ NameError trong bộ xử lý pha LAND:} Khắc phục lỗi thiếu biến thời gian \texttt{now} trong hàm xử lý hạ cánh của \texttt{mission\_fsm\_node.py}, đảm bảo luồng điều khiển vận tốc luôn liên tục cho tới khi drone nằm yên trên sàn.
    \item \textbf{Tối ưu hóa tài nguyên và dọn dẹp tiến trình:} Triển khai cơ chế quản lý theo nhóm tiến trình (\texttt{os.killpg} và \texttt{clean\_simulation\_processes}), chấm dứt 100\% tiến trình ROS 2 và Gazebo con sau mỗi chuyến bay, giữ bộ nhớ RAM luôn trống $> 7.5\text{ GB}$ và triệt tiêu hoàn toàn giật lag.
\end{enumerate}

% ==============================================================================
% PHẦN VI: PHÂN TÍCH GIÁ THÀNH & TÍNH KHẢ THI KINH TẾ (BUSINESS CASE)
% ==============================================================================
\section{PHÂN TÍCH GIÁ THÀNH VÀ HIỆU QUẢ KINH TẾ SO VỚI THỊ TRƯỜNG}

\subsection{Bảng So sánh Đa chiều với các Giải pháp trên Thị trường}

\begin{table}[h!]
\centering
\resizebox{\textwidth}{!}{%
\begin{tabular}{@{}lp{2.6cm}p{2.6cm}p{2.6cm}p{2.6cm}p{3.0cm}@{}}
\toprule
\textbf{Tiêu chí} & \textbf{RTK-GPS Công nghiệp} & \textbf{Mạng UWB / Beacon} & \textbf{3D LiDAR SLAM} & \textbf{Drone Dock Thương mại} & \textbf{Giải pháp Sáng kiến VDT} \\ \midrule
\textbf{Cơ chế hoạt động} & Base RTK cố định + Rover RTK trên drone & 4-6 trạm phát UWB quanh bãi + Tag trên drone & Cảm biến LiDAR quét mây điểm + Máy tính GPU & Trạm cơ khí khép kín + Drone chuyên dụng độc quyền & \textbf{Thị giác mép Pi 5 + Mốc kép A0 + Dẫn đường SMC} \\ \midrule
\textbf{Chi phí phần cứng bổ sung (BOM)} & \textbf{\$1,500 -- \$4,500}\newline(38 -- 115 tr VNĐ) & \textbf{\$800 -- \$2,200}\newline(20 -- 55 tr VNĐ) & \textbf{\$3,500 -- \$10,000}\newline(90 -- 250 tr VNĐ) & \textbf{\$15,000 -- \$35,000}\newline(380 -- 880 tr VNĐ) & \textbf{\textasciitilde \$380 (\textasciitilde 9.5 triệu VNĐ)}\newline\textit{(Tiết kiệm 85\% -- 95\%)} \\ \midrule
\textbf{Hạ tầng bãi đáp} & Trạm Base đo đạc tọa độ cố định, cấp điện & 4 cột mốc cơ khí cần cấp pin, dễ lệch & Cần bãi có đặc trưng phản xạ LiDAR & Trạm dock cơ khí nặng 50-100kg, cấp điện lưới & \textbf{Tấm bạt A0 (\$5 \textasciitilde 120k), không cần điện, dán trong 10s} \\ \midrule
\textbf{Sai số tiếp đất ($R_{xy}$)} & $2 - 5\text{ cm}$ (lý tưởng) & $10 - 20\text{ cm}$ (nhiễu phản xạ) & $3 - 8\text{ cm}$ & $2 - 5\text{ cm}$ (dùng ngàm ép) & \textbf{$1.0 - 2.0\text{ cm}$ (Đã kiểm chứng)} \\ \midrule
\textbf{Mất sóng GPS (GPS-denied)} & \textcolor{alertred}{\textbf{Liệt hoàn toàn}} & Vùng phủ hẹp ($< 20\text{m}$) & Tốt trong nhà & Phụ thuộc GPS trạm & \textbf{Tự hành 100\% nhờ Marker-Relative} \\ \midrule
\textbf{Bám bãi di động} & Cực kỳ phức tạp & Khó căn chỉnh & Rất nặng thuật toán & \textcolor{alertred}{\textbf{Không hỗ trợ}} & \textbf{Tự bám xe/thuyền $\le 1\text{m/s}$ nhờ I-APF} \\ \midrule
\textbf{Mức làm chủ công nghệ} & Phụ thuộc hãng & Hộp đen driver & Nặng phần cứng GPU & Đóng mã hoàn toàn & \textbf{Làm chủ 100\% mã nguồn C++} \\ \bottomrule
\end{tabular}%
}
\caption{So sánh toàn diện giữa sáng kiến VDT và các giải pháp hiện có trên thị trường}
\label{tab:market_comparison}
\end{table}

\subsection{Bóc tách Chi phí Bill of Materials (BOM) Gói Nâng cấp Thông minh}

\begin{table}[h!]
\centering
\small
\begin{tabularx}{\textwidth}{@{}l X r r@{}}
\toprule
\textbf{Hạng mục linh kiện} & \textbf{Thông số kỹ thuật lựa chọn} & \textbf{Đơn giá (VNĐ)} & \textbf{Tỷ trọng (\%)} \\ \midrule
Máy tính nhúng Companion & Raspberry Pi 5 (4GB RAM) & 1.950.000 & 20.8\% \\
Cảm biến thị giác độ sâu & Intel RealSense D435i / D430 & 7.200.000 & 77.0\% \\
Cơ cấu Gimbal Pitch & 1× Micro Servo MG90S + Khung in 3D & 80.000 & 0.9\% \\
Bãi đáp mục tiêu mốc kép & Decal/Bạt in chuẩn A0 Nested ArUco & 120.000 & 1.3\% \\ \midrule
\textbf{TỔNG CHI PHÍ GÓI THÔNG MINH} & \textbf{Gói phụ trợ (Smart Add-on Kit)} & \textbf{9.350.000 VNĐ} & \textbf{100.0\%} \\ \bottomrule
\end{tabularx}
\caption{Bóc tách chi phí linh kiện thực tế của giải pháp đề xuất}
\label{tab:bom_cost}
\end{table}

\subsection{Hiệu quả Kinh tế và Giá trị Gia tăng đối với Doanh nghiệp}
\begin{enumerate}[noitemsep,topsep=2pt]
    \item \textbf{Khả năng nhân rộng đại trà (High Scalability):} Thay vì phải đầu tư hàng trăm triệu đồng cho mỗi điểm hạ cánh, doanh nghiệp chỉ cần chi chưa đầy $10$ triệu đồng để nâng cấp bất kỳ máy bay không người lái tiêu chuẩn nào hiện có thành hệ thống tự hành thông minh.
    \item \textbf{Cơ động dã chiến không giới hạn:} Bãi đáp chỉ là một tấm bạt mỏng có thể gấp gọn trong balo. Trong các nhiệm vụ cứu hộ, trinh sát hoặc kiểm tra công nghiệp, bãi đáp có thể trải ra nóc xe bán tải hoặc boong tàu chỉ trong 10 giây.
    \item \textbf{Tự chủ công nghệ và An toàn thông tin:} Toàn bộ thuật toán được xây dựng trên nền tảng mã nguồn mở chuẩn công nghiệp (ROS 2 Humble, PX4 Autopilot, C++20), giúp doanh nghiệp hoàn toàn làm chủ hệ thống, dễ dàng tùy biến giao thức bảo mật theo tiêu chuẩn riêng của công ty mà không lo phụ thuộc hay rò rỉ dữ liệu sang nhà cung cấp thứ ba.
\end{enumerate}

% ==============================================================================
% PHẦN VII: KẾT LUẬN & LỘ TRÌNH TRIỂN KHAI (ROADMAP)
% ==============================================================================
\section{KẾT LUẬN VÀ LỘ TRÌNH TRIỂN KHAI TIẾP THEO}

\subsection{Kết luận}
Sáng kiến \textbf{"Hệ thống Quadrotor tự hành bám mục tiêu di động, né vật cản 3D và hạ cánh chính xác chi phí thấp"} đã chứng minh tính khả thi vượt bậc về cả ba phương diện:
\begin{itemize}[noitemsep,topsep=2pt]
    \item \textbf{Học thuật \& Thuật toán:} Giải quyết trọn vẹn các thách thức phi tuyến phức tạp (bẫy cực tiểu địa phương, hiện tượng GNRON, góc nhìn camera khi hạ dốc, và kẹt lơ lửng).
    \item \textbf{Kỹ thuật Nhúng:} Chuyển thể thành công sang mã nguồn C++20 Header-Only chạy mượt mà trên phần cứng hạn chế tài nguyên Raspberry Pi 5 ($< 2\text{ms}$/chu trình).
    \item \textbf{Hiệu quả Kinh tế:} Đạt độ chính xác tiếp đất cấp độ centimet ($1 - 2\text{cm}$) với chi phí chỉ bằng $5\% - 15\%$ so với các thiết bị công nghiệp ngoại nhập.
\end{itemize}

\subsection{Lộ trình Phát triển và Kiến nghị (Roadmap)}
\begin{itemize}[noitemsep,topsep=2pt]
    \item \textbf{Giai đoạn 1 (Đã hoàn thành 100\%):} Thiết kế thuật toán, kiểm chứng mô phỏng Gazebo Harmonic / PX4 SITL, đóng gói C++20 nhúng trên Raspberry Pi 5.
    \item \textbf{Giai đoạn 2 (Tháng 10 -- 11/2026):} Thực hiện bay thử nghiệm ngoài trời trên khung Holybro S500 thực tế, hiệu chỉnh tham số điều khiển theo điều kiện gió thực tế.
    \item \textbf{Giai đoạn 3 (Quý 1/2027):} Tích hợp các cực sạc tiếp xúc tự động trên bề mặt bãi đáp A0, hoàn thiện giải pháp trạm bay tự hành không người trực (Autonomous Drone-in-a-Box) thương mại hoàn chỉnh với tổng giá thành dưới 50 triệu VNĐ.
\end{itemize}

\vspace{1.5em}
\noindent
\textbf{Kiến nghị:} Nhóm kính đề nghị Hội đồng Giám đốc công ty phê duyệt sáng kiến ý tưởng, hỗ trợ kinh phí phục vụ đợt bay thử nghiệm ngoài trời và xem xét bảo hộ giải pháp kỹ thuật cấp doanh nghiệp.

\vspace{2.5em}
\begin{flushright}
    \textbf{Đại diện Nhóm Tác giả Sáng kiến}\\[0.3em]
    \textit{(Ký và ghi rõ họ tên)}\\[3.5em]
    \textbf{Kỹ sư Duy (Lead R\&D / Algorithm)}
\end{flushright}

\end{document}
